\documentclass[10pt,a4paper]{amsart}
\usepackage[margin=19mm]{geometry}
\usepackage{amsmath,amssymb,mathtools}
\usepackage[scr=rsfs,cal=euler]{mathalfa}
\usepackage{xfrac}
\usepackage[unicode,hidelinks,bookmarks=false]{hyperref}
\newcommand{\ie}{i.e.\ }
\newcommand{\cf}{cf.\ }
\newcommand{\resp}{respectively\ }
\newcommand{\Subsection}{\subsection}
\providecommand{\nobreakdash}{\nobreak\hbox{-}}
\setlength{\emergencystretch}{2em}
\multlinegap0pt
\usepackage{bm}
\usepackage{bbm}
\usepackage{dsfont}
\usepackage{xspace}
\usepackage{cancel}
\usepackage{mathtools}
\usepackage{amsthm}
\makeatletter
\@ifundefined{theorem}{\newtheorem{theorem}{Theorem}}{}
\@ifundefined{lemma}{\newtheorem{lemma}[theorem]{Lemma}}{}
\makeatother
\newcommand\supp{\operatorname{supp}}
\newcommand{\ff}{\varphi}
\newcommand{\intt}{\mbox{\rm int\,}}
\title[Duality for Delsarte's extremal problem on compact Gelfand p]{Duality for Delsarte's extremal problem on compact Gelfand pairs}
\author{Elena E. Berdysheva, B\'alint Farkas, Marcell Ga\'al, Mita D. Ramabulana, Szil\'ard Gy. R\'ev\'esz}
\date{Source: arXiv:2603.11792v4 (CC BY 4.0)}
\begin{document}
\maketitle

\noindent\emph{Source and licence.} Duality for Delsarte's extremal problem on compact Gelfand pairs. Version arXiv:2603.11792v4 (CC BY 4.0), distributed under the Creative Commons Attribution 4.0 International licence, as recorded in the arXiv licence metadata for this exact version and shown on the abstract page https://arxiv.org/abs/2603.11792v4. Full rights notice: licenses/arxiv-2603.11792-CC-BY-4.0.txt.\par\smallskip

\noindent\textit{Editorial scope.} Counted unit: the Theorem whose source label is \texttt{main\_theorem} in the article (source0.tex lines 1052--1059, source label \texttt{main\_theorem}), delivered together with the article's own complete original proof of it (source0.tex lines 1061--1105, one \texttt{proof} environment of 45 lines). No mathematics is written, repaired or re-derived: statement and proof are the article's own text line for line. Required context retained uncounted: eq:C\_correctly\_posed (lines 833--835); lemma:C\_epsilon (lines 904--913). Every definition, display and label of those lines is retained unchanged. Omitted from this package and not redistributed: 1--832, 836--903, 914--1051, 1060--1060, 1106--1446. Nothing inside a retained excerpt is omitted or altered: every retained body is delivered exactly as the article's lines are written, so no omission receipt of any kind accompanies this package. Citations: the retained excerpts carry no citation command at all, so no bibliography entry of the article is transcribed and no reference is redistributed. Reference adaptation: none was needed and none was made; every reference inside the delivered unit resolves inside it, so no unresolved reference or citation is produced. Printed slips kept literal: no printed word, symbol, hypothesis or number of the counted statement or of its proof was altered. Layout and class: the article's own class is \texttt{amsart} with packages including inputenc, amsfonts, amsmath, amsthm, amssymb, bbm, verbatim, mathabx, xcolor, comment, color, ulem, babel; the wrapper is the lane's pinned \texttt{amsart} editorial wrapper at 10pt on A4, which restates the theorem-like environments the retained text needs and adds the article's own macro definitions that the retained text needs (\texttt{\textbackslash ff}, \texttt{\textbackslash intt}, \texttt{\textbackslash supp}), with extra wrapper packages (bm, bbm, dsfont, xspace, cancel, mathtools). The delivered numbering is the unit's own. Assets: no retained span contains a figure environment, a graphics-inclusion command, a PDF-page inclusion, a table or a TikZ picture; no image file is copied into this package, the package declares no asset dependency and no image file is redistributed with it. Licence: version arXiv:2603.11792v4 of this article is distributed under the Creative Commons Attribution 4.0 International licence, as recorded in the arXiv licence metadata for this exact version and shown on the abstract page https://arxiv.org/abs/2603.11792v4. A full rights notice is delivered at licenses/arxiv-2603.11792-CC-BY-4.0.txt. Characters: every retained excerpt is pure ASCII, so the delivered engine, XeLaTeX with its default Latin Modern text font, reports no missing character.
\begin{equation} \label{eq:C_correctly_posed}
\lim_{\varepsilon \to 0+} u_\Gamma^{\sigma}(\Omega_+, \Omega_-;\varepsilon) = u_\Gamma^{\sigma}(\Omega_+, \Omega_-).
\end{equation}

\begin{lemma} \label{lemma:C_epsilon}

Let $(G,K)$ be a compact Gelfand pair. Let $\Omega_+$, $\Omega_-$ be $K$-bi-invariant symmetric Borel subsets of $G$ with $e \in \intt{\Omega_+}$, and  $\Omega_+^c$, $\Omega_-^c$ satisfying Assumption~O. 
Assume that $\sigma\in M(G)^K$ has the form $\sigma = \delta_e^K + \tau $, where \(\tau \in M(G)^K\) is positive definite and absolutely continuous with respect to the Haar measure $\lambda_G$.
Then
$$
\lim_{\varepsilon \to 0+} \mathcal{A}^{K,\sigma}_G(\Omega_+, \Omega_-;\varepsilon) = \mathcal{A}^{K,\sigma}_G(\Omega_+, \Omega_{-}).
$$

\end{lemma}

\begin{theorem}\label{main_theorem}

Let $(G,K)$ be a compact Gelfand pair. Let $\Omega_+$, $\Omega_-$ be $K$-bi-invariant symmetric Borel subsets of $G$ with $e \in \intt{\Omega_+}$, and $\Omega_+^c$, $\Omega_-^c$ satisfying Assumption O. Assume that $\sigma\in M(G)^K$ has the form $\sigma = \delta_e^K + \tau $, where \(\tau \in M(G)^K\) is positive definite and absolutely continuous with respect to the Haar measure $\lambda_G$. Then
$$
u_\Gamma^{\sigma}(\Omega_+, \Omega_-) = v_\Gamma^{\sigma}(\Omega_+, \Omega_-).
$$

\end{theorem}

\begin{proof}

To prove the theorem, we must show \eqref{eq:C_correctly_posed}.

The Delsarte-type problem equivalent to $u_\Gamma^{\sigma}(\Omega_+, \Omega_-;\varepsilon)$ differs from $\mathcal{A}^{K,\sigma}_G(\Omega_+, \Omega_-;\varepsilon)$ by the normalisation: instead of $\langle \ff, \sigma \rangle_2 = 1$, we need the normalisation $\int_G \ff \,\mbox{d}\lambda_G =1$. Therefore, for the sake of this proof we introduce the extremal constant, with the notation $s = \widehat{\sigma}$,
\begin{align*}
\mathcal{A}_G^{K,\sigma,1}(\Omega_+, \Omega_-; \varepsilon) & := 
\sup \left\{ \frac{\int_G \ff \,\mbox{d} \lambda_G}{\langle \ff, \sigma \rangle_2} : \ff \in C(G)^K, \ \ff \gg 0, \right.\\
& \left.
\hspace{14mm} \int_G \ff \,\mbox{d}\lambda_G =1, \ \ff|_{\Omega_+^c} \le \varepsilon, \ \ff|_{\Omega_-^c} \ge -\varepsilon
\right\} 
 \\
& = \sup \left\{ \frac{ 1 }{ s(\mathbbm{1}_G) + \sum_{\gamma \in \supp{f} \setminus \{\mathbbm{1}_G\}} f(\gamma) s(\gamma) }  : \right.   f \in \ell^1_+(\Gamma), \
\varphi|_A \le \varepsilon, \\ 
& \hspace{14mm} \varphi|_B \ge -\varepsilon, \text{ where }\left. \varphi = \mathbbm{1}_G + \sum_{\gamma \in \supp{f} \setminus \{\mathbbm{1}_G\}} f(\gamma) \gamma
\right\}.
\end{align*}
We have
$$
\mathcal{A}_G^{K,\sigma,1}(\Omega_+, \Omega_-;\varepsilon) = \frac{1}{s(\mathbbm{1}_G) + u_\Gamma^{\sigma}(\Omega_+, \Omega_-;\varepsilon)}.
$$
Note that admissible functions $\varphi$ in the problem $\mathcal{A}_G^{K,\sigma,1}(\Omega_+, \Omega_-;\varepsilon)$ satisfy $\int_G \varphi \,\mbox{d}\lambda_G = 1$. Since  $\int_G \varphi \,\mbox{d}\lambda_G \le \varphi(e) \lambda_G(G) = \varphi(e)$, we also have $\varphi(e) \ge 1$, and therefore $\langle \ff, \sigma \rangle_2 \ge 1$. Now take $\psi := \frac{\varphi}{\langle \ff, \sigma \rangle_2}$. We have $\psi \in C(G)^K$, $\psi \gg 0$, $\langle \psi, \sigma \rangle_2 = 1$, $\psi|_A \le \frac{\varepsilon}{\langle \ff, \sigma \rangle_2} \le \varepsilon$, $\psi|_B \ge -\frac{\varepsilon}{\langle \ff, \sigma \rangle_2} \ge -\varepsilon$. 
Thus, $\psi \in \mathcal{F}^{K,\sigma}_G(\Omega_+, \Omega_-; \varepsilon)$, and
$$
\frac{\int_G \ff \, \mbox{d}\lambda_G}{\langle \ff, \sigma \rangle_2} = \int_G \psi \, \mbox{d}\lambda_G.
$$
Hence,
$$
\mathcal{A}_G^{K,\sigma,1}(\Omega_+, \Omega_-;\varepsilon) \le \mathcal{A}^{K,\sigma}_G(\Omega_+, \Omega_-;\varepsilon).
$$
Since $u_\Gamma^{\sigma}(\Omega_+, \Omega_-;\varepsilon)$ is an extension of the minimization problem $u_\Gamma^\sigma(\Omega_+,\Omega_-)$, we have that $u_\Gamma^{\sigma}(\Omega_+, \Omega_-) \ge u_\Gamma^{\sigma}(\Omega_+, \Omega_-;\varepsilon)$. Thus,
\begin{align*}
\mathcal{A}^{K,\sigma}_G(\Omega_+, \Omega_{-}) & = \frac{1}{s(\mathbbm{1}_G) + u_\Gamma^{\sigma}(\Omega_+, \Omega_-)} 
\le \frac{1}{s(\mathbbm{1}_G) + u_\Gamma^{\sigma}(\Omega_+, \Omega_-;\varepsilon)} \\
& = \mathcal{A}_G^{K,\sigma,1}(\Omega_+, \Omega_-;\varepsilon) \le \mathcal{A}^{K,\sigma}_G(\Omega_+, \Omega_-;\varepsilon).
\end{align*}
By Lemma~\ref{lemma:C_epsilon},
$\lim_{\varepsilon \to 0+} \mathcal{A}^{K,\sigma}_G(\Omega_+, \Omega_-;\varepsilon) = \mathcal{A}^{K,\sigma}_G(\Omega_+, \Omega_{-})$,
which in turn implies~\eqref{eq:C_correctly_posed}. We conclude the problem $u_\Gamma^{\sigma}(\Omega_+, \Omega_-)$ is well-posed. This implies that the strong duality relation
$$
u_\Gamma^{\sigma}(\Omega_+, \Omega_-) = v_\Gamma^{\sigma}(\Omega_+, \Omega_-)
$$
holds true.

\end{proof}

\end{document}
